%%%%%%%%%%%%%%%%%%%%%%%%%%%%%%%%%%%%%%%%%%%%%%%%%%%%%%%%%%%%%%%%%%%%%%%%%
% If you want to use arara, you need the following directives:          %
% In the preamble:                                                      %
% arara: pdflatex: { synctex: yes, shell: yes, interaction: batchmode }
% arara: pdflatex: { synctex: yes, shell: yes, interaction: batchmode }
% arara: bibtex
% arara: pdflatex: { synctex: yes, shell: yes, interaction: batchmode }
% arara: clean: { extensions: [ aux, bbl, bcf, blg, glg, glstex, lof, log, ptc, run.xml, thm, toc ] }
%%%%%%%%%%%%%%%%%%%%%%%%%%%%%%%%%%%%%%%%%%%%%%%%%%%%%%%%%%%%%%%%%%%%%%%%%

% !TEX TS-program = pdflatex
% !TeX encoding = UTF-8
% !TeX spellcheck = eng_GB

%%%% ---- LaTeX File ---- %%%%
\providecommand{\PathP}{./Parts/} % Path to Parts folders
\providecommand{\PathR}{./Resources/} % Path to Resources folders
\providecommand{\PathSO}{./Sources/} % Path to Sources folders
\providecommand{\PathS}{./Supplies/} % Path to Supplies folders
\providecommand{\lang}{english}	% english, french
\providecommand{\BuildTrigger}{buildnew} % Trigger to build or not standalone parts.

\makeatletter
	\ifx\c@levelstanda\undefined % to create a new-counter if it doesn't exists.
		\newcounter{levelstanda}
		\def\input@path{ {\PathSO}{./another-directory/} } % to be able to find package in source folders.
	\else
		{}% it already exist
	\fi
\makeatother


%%%%---- DOCUMENT-CLASS -------------------------%%%%
\documentclass[a4paper]{article}  % or another class
	\usepackage[l2tabu, orthodox]{nag}	% Helps to correct the script

%%%%---- OS & FONT -------------------------%%%%
	\usepackage[utf8]{inputenc} % Input keyboard encoding: latin1=windows, utf8=linux, applemac=Mac
	\usepackage[T1]{fontenc} 	% Output font encoding: T1 for Europeans languages.

%%%%---- PERSONAL Custom ----------------------%%%%
	\newcommand{\Pathbasebib}{\PathR/ref_RN.bib,\PathR/ref_XX.bib}% Careful no space!
	\newcommand{\PathSymbolsEntries}{\PathR/UsersSymbolsAcronyms.bib}
	\input{\PathR/UsersPackagesDefs.tex}
	\input{\PathR/UsersMathsCommands.tex}
		\geometry{margin=2.5cm}

%%%%%---- JOURNAL SPECS -------------------------%%%%
	%\journal{Journal of XX}
	\title{My Empty Template}
	\author{Romain \famName{Noël}}
	\date{\today}


%%%%%%%%%%%%%%%  MAIN DOCUMENT  %%%%%%%%%%%%%%%%%%%%%
\begin{document}
\stepcounter{levelstanda}

%% FRONTMATTER
\maketitle

	\begin{abstract}
		This document is prepared using LaTeX2e~\cite{Lamport94} and shows the desired format and appearance of a manuscript prepared for the Proceedings of the SPIE.\footnote{The basic format was developed in 1995 by Rick Hermann (SPIE) and Ken Hanson (Los Alamos National Lab.).}
		It contains general formatting instructions and hints about how to use LaTeX.
		The LaTeX source file that produced this document, {\ttfamily article.tex} (Version 3.4), provides a template, used in conjunction with {\ttfamily spie.cls} (Version 3.4).
		These files are available on the Internet at \url{https://www.overleaf.com}.
		The font used throughout is the LaTeX default font, Computer Modern Roman, which is equivalent to the Times Roman font available on many systems.
	\end{abstract}

%% MAINMATTER
\input{\PathP/example.tex}
\input{\PathP/new_part.tex}


%% BACK-MATTER
%\appendix
\ifnum\value{levelstanda}>1
	{}% Not in the main.
\else
	% References
	\bibliography{\Pathbasebib} % for bibtex or natbib
	%\printbibliography     % for biblatex

	%%% To comment for the final version
	%\tableofcontents
	\printunsrtglossary[type=acronym]
	\printunsrtglossary[type=symbol]
\fi

\addtocounter{levelstanda}{-1}
\end{document}
%EoF